\documentclass[11pt,a4paper]{article}
\usepackage[utf8]{inputenc}
\usepackage[T1]{fontenc}
\usepackage{amsmath,amssymb,amsthm}
\usepackage[margin=1in]{geometry}
\usepackage{enumitem}

\theoremstyle{plain}
\newtheorem{theorem}{Theorem}
\newtheorem{corollary}[theorem]{Corollary}
\newtheorem{proposition}[theorem]{Proposition}
\theoremstyle{definition}
\newtheorem{definition}[theorem]{Definition}
\newtheorem{assumption}[theorem]{Assumption}
\theoremstyle{remark}
\newtheorem{remark}[theorem]{Remark}

\title{Reading Report: Mayer--Vietoris and Twisted \v{C}ech Spectral Sequences\\
for $C^*$-Algebras with Free Quantum Group Coefficients}
\author{Auto-generated report \\ \small arXiv:2602.22708 --- Takao Inou\'e (2026)}
\date{}

\begin{document}
\maketitle

%------------------------------------------------------------
\section{Core Question}
%------------------------------------------------------------
How can one systematically detect \emph{higher} (i.e.\ beyond the classical six-term exact sequence) $K$-theoretic obstructions in $C^*$-algebras assembled by gluing local pieces, when the fibers carry the additional structure of a free quantum group crossed product?
In particular, when does a ``twist'' in the gluing data produce a $C^*$-algebra that is \emph{not} Morita equivalent to the trivial (untwisted) field?

%------------------------------------------------------------
\section{Main Statement}
%------------------------------------------------------------
\begin{theorem}[Identification of $d_2$, Theorem 4.4 in the paper]
Let $X = S^2$ equipped with a good open cover $\mathcal{U}$, and let
\[
D \;:=\; \mathcal{O}_n \rtimes_r \widehat{O_n^+}
\]
be the reduced crossed product of the Cuntz algebra $\mathcal{O}_n$ by the discrete dual $\widehat{O_n^+}$ of the free orthogonal quantum group $O_n^+$.
Fix:
\begin{itemize}[nosep]
  \item a nontrivial \v{C}ech $2$-cocycle $c \in \check{H}^2(S^2;\mathbb{Z}/2) \cong \mathbb{Z}/2$,
  \item an involutive automorphism $\phi \in \mathrm{Aut}(D)$ with $\phi^2 = \mathrm{id}$.
\end{itemize}
Let $A_c := C^*(G_{\mathcal{U}}, B_{\phi,c})$ be the twisted Fell-bundle $C^*$-algebra over the \v{C}ech groupoid $G_{\mathcal{U}}$.
Under a Kirchberg--UCT hypothesis on $D$, the \v{C}ech--Mayer--Vietoris spectral sequence satisfies
\[
  d_2 \;:\; E^2_{2,q} \;\longrightarrow\; E^2_{0,q+1}, \qquad
  d_2 \;=\; \phi_* - \mathrm{id} \;:\; K_q(D) \to K_{q+1}(D).
\]
\end{theorem}

\begin{corollary}[Morita obstruction]
If $K_0(D) \cong K_1(D) \cong \mathbb{Z}/3$ and $\phi_* = -1$, then $d_2 = -2 \equiv 1 \pmod{3}$ is an isomorphism, so
\[
  K_*(A_c) \;\not\cong\; K_*(C(S^2) \otimes D),
\]
and the twisted algebra $A_c$ is \textbf{not} Morita equivalent to the trivial field $A_0 = C(S^2) \otimes D$.
\end{corollary}

%------------------------------------------------------------
\section{A Concrete Example First}
%------------------------------------------------------------
Take $n = 4$, so $D = \mathcal{O}_4 \rtimes_r \widehat{O_4^+}$.
Suppose (via a quantum Pimsner--Voiculescu computation) that
$K_0(D) \cong K_1(D) \cong \mathbb{Z}/3$ and there exists
$\phi \in \mathrm{Aut}(D)$ with $\phi_* = -\mathrm{id}$ on $K$-groups.

\medskip\noindent
\textbf{Untwisted algebra} $A_0 = C(S^2) \otimes D$:\;
By K\"unneth, $K_*(A_0) \cong K_*(D) \oplus K_*(D) \cong (\mathbb{Z}/3)^{\oplus 2}$ in each degree.

\medskip\noindent
\textbf{Twisted algebra} $A_c$:\;
The spectral sequence has $E^2_{2,q} \cong E^2_{0,q+1} \cong \mathbb{Z}/3$.
The differential $d_2 = \phi_* - \mathrm{id} = -1 - 1 = -2 \equiv 1 \pmod{3}$
is an isomorphism, killing the $p = 2$ contribution at $E^\infty$.
Hence $K_*(A_c) \cong \mathbb{Z}/3$ (one copy, not two), so $A_c \not\sim_M A_0$.

%------------------------------------------------------------
\section{Intuition and Setup}
%------------------------------------------------------------

\subsection*{Key objects and notation}
\begin{itemize}[nosep]
  \item $\mathcal{O}_n$: the \textbf{Cuntz algebra}, the universal $C^*$-algebra generated by $n$ isometries $S_1, \ldots, S_n$ satisfying $\sum_{i=1}^n S_i S_i^* = 1$.
  \item $O_n^+$: the \textbf{free orthogonal quantum group} (Wang, 1995), a compact quantum group whose representation ring generalizes that of $O(n)$.  Its discrete dual $\widehat{O_n^+}$ acts on $\mathcal{O}_n$ via a coaction $\delta(S_i) = \sum_j S_j \otimes u_{ji}$.
  \item $D = \mathcal{O}_n \rtimes_r \widehat{O_n^+}$: the \textbf{reduced crossed product}, packaging the quantum symmetry into a single ``coefficient algebra''.
  \item \textbf{Equivariant ideal cover}: given $A = \sum I_i$ (sum of closed two-sided ideals), the short exact sequence
    \[
      0 \to I_1 \cap I_2 \xrightarrow{\Delta} I_1 \oplus I_2 \xrightarrow{\Sigma} A \to 0
    \]
    produces the classical \textbf{Mayer--Vietoris (MV) six-term exact sequence} in $K$-theory.
  \item \textbf{\v{C}ech spectral sequence}: for a finite cover $\{I_i\}$, one forms the chain complex
    \[
      C^{(q)}_p = \bigoplus_{i_0 < \cdots < i_p} K_q(I_{i_0 \cdots i_p}), \qquad d_1 = \text{alternating restriction}.
    \]
    The $E^2$-page is the \v{C}ech homology of the nerve with $K$-theory coefficients.
  \item \textbf{Kirchberg--UCT hypothesis}: $D$ is separable, nuclear, simple, purely infinite, and satisfies the Universal Coefficient Theorem.  This ensures that $K$-theory classifies $D$ up to isomorphism (Kirchberg--Phillips).
\end{itemize}

\subsection*{Big picture}
The MV six-term sequence is a well-known tool for computing $K$-theory of $C^*$-algebras built by gluing.  It detects obstructions at the level of \emph{one-dimensional} overlaps ($d_1$).
However, when gluing data carries \emph{higher cohomological} information (here a $\mathbb{Z}/2$-valued \v{C}ech 2-cocycle on $S^2$), the six-term sequence is blind to it.
The paper upgrades the MV sequence to a full \v{C}ech spectral sequence, where the \emph{second differential} $d_2$ captures this extra twist---and shows it can obstruct Morita equivalence.

%------------------------------------------------------------
\section{Main Results (Bulleted)}
%------------------------------------------------------------
\begin{enumerate}[label=\textbf{(\roman*)}]
  \item \textbf{MV as \v{C}ech $d_1$:}
    The classical MV boundary map $\partial$ is identified with the \v{C}ech differential $d_1$ (alternating restriction) in a spectral sequence whose $E^1$-page is built from iterated intersections of an equivariant ideal cover.

  \item \textbf{Coefficient reduction:}
    When the quantum group acts trivially on the base $X$ and all intersections $U_{i_0 \cdots i_p}$ are contractible,
    \[
      K_q(I_{i_0 \cdots i_p}) \cong K_{q-1}(D),
    \]
    so the entire $E^1$-page is determined by $K_*(D)$ and the combinatorics of the nerve of the cover.

  \item \textbf{Degeneration for contractible bases:}
    For $X = [0,1]$ with a good cover, the spectral sequence degenerates at $E^2$, recovering $K_q(A) \cong K_{q-1}(D)$.

  \item \textbf{Twisted $d_2$ formula:}
    For $X = S^2$ with a nontrivial $\mathbb{Z}/2$-cocycle $c$ and involution $\phi$,
    \[
      d_2 = \phi_* - \mathrm{id}.
    \]

  \item \textbf{Morita obstruction:}
    In the cyclic regime ($K_*(D) \cong \mathbb{Z}/3$, $\phi_* = -1$), $d_2$ is an isomorphism and $A_c$ is not Morita equivalent to $C(S^2) \otimes D$.
\end{enumerate}

%------------------------------------------------------------
\section{Proof Sketch}
%------------------------------------------------------------

\begin{enumerate}[label=\textbf{Step \arabic*.}, leftmargin=2.5em]
  \item \textbf{Set up the equivariant ideal cover.}
    Start with $B = C(X) \otimes \mathcal{O}_n$ carrying the coaction $\alpha = \mathrm{id}_{C(X)} \otimes \delta$, and set $A = B \rtimes_r \widehat{O_n^+}$.
    An open cover $\{U_i\}$ of $X$ gives equivariant ideals $I_i = C_0(U_i) \otimes D$.

  \item \textbf{Build the \v{C}ech spectral sequence.}
    Form the chain complex $C^{(q)}_p = \bigoplus K_q(I_{i_0 \cdots i_p})$ with $d_1$ = alternating restriction.
    Under the contractibility hypothesis on intersections, apply the coefficient reduction: all entries become copies of $K_{q-1}(D)$.

  \item \textbf{Identify $E^2$ for $X = S^2$.}
    The nerve of a good cover of $S^2$ has the homology of $S^2$, so $E^2_{p,q}$ is nonzero only for $p = 0$ and $p = 2$.
    Specifically, $E^2_{0,q} \cong K_q(D)$ and $E^2_{2,q} \cong K_q(D)$.

  \item \textbf{Introduce the twist.}
    Replace the trivial transition data with a $\mathbb{Z}/2$-valued \v{C}ech 2-cocycle $c$ (nontrivial class in $\check{H}^2(S^2;\mathbb{Z}/2)$) and an involution $\phi \in \mathrm{Aut}(D)$.
    Realize this as a twisted Fell bundle over the \v{C}ech groupoid $G_{\mathcal{U}}$.

  \item \textbf{Compute $d_2$.}
    The twist is invisible to $d_1$ (which depends only on double overlaps).
    On triple overlaps, traversing a $2$-simplex in the nerve introduces the automorphism $\phi$ (controlled by $c_{ijk}$).
    The resulting compatibility defect on $K$-theory is exactly $\phi_* - \mathrm{id}$, yielding $d_2 = \phi_* - \mathrm{id}$.

  \item \textbf{Specialize to the cyclic regime.}
    With $K_*(D) \cong \mathbb{Z}/3$ and $\phi_* = -1$:
    $d_2 = -1 - 1 = -2 \equiv 1 \pmod{3}$, an isomorphism on $\mathbb{Z}/3$.
    The $p = 2$ column is killed at $E^3 = E^\infty$, so
    $K_*(A_c) \cong \mathbb{Z}/3 \not\cong (\mathbb{Z}/3)^2 \cong K_*(A_0)$.

  \item \textbf{Conclude Morita non-equivalence.}
    Since Morita equivalence preserves $K$-theory, $A_c \not\sim_M A_0$. \qed
\end{enumerate}

\end{document}
